%% ma: L12-B02-0017
%% lop: 12
%% chuong: 1
%% bai: 2
%% dang: 12.2.10
%% loai: TN4
%% muc: TH
%% dap_an: "B"
%% nguon: "Ảnh giáo viên gửi 05/11/2026 (M11): Kiểm tra giữa HK1 Toán 12 (2026-2027), THPT Hồ Thị Bi (TP.HCM), mã đề 1133, Phần I Câu 7"
%% kien_thuc: "Bài 2 (giá trị lớn nhất trên nửa khoảng dùng tính đơn điệu); Bài 1 (đạo hàm của thương)"
%% trang_thai: chinh-thuc
%% ghi_chu: "Đề gốc không kèm đáp án; đã giải độc lập. Hỏi trên nửa khoảng (dạng 12.2.10 gốc trên đoạn), cùng cách làm."
\begin{ex}
Giá trị lớn nhất của hàm số $f(x)=\dfrac{x+4}{x-1}$ trên nửa khoảng $[5;+\infty)$ bằng
\choice{$\dfrac{-5}9$}{\True $\dfrac94$}{$\dfrac{-5}{16}$}{$\dfrac{16}5$}
\loigiai{$f'(x)=\dfrac{(x-1)-(x+4)}{(x-1)^2}=\dfrac{-5}{(x-1)^2}<0$ với mọi $x\ge5$ nên hàm số nghịch biến trên $[5;+\infty)$.
Vậy $\max\limits_{[5;+\infty)}f(x)=f(5)=\dfrac94$.}
\end{ex}
